\section{Klassendiagramm-Notation}


\subsection{Beziehungen}
\begin{itemize}
    \item \textbf{Vererbung:} \texttt{Drucksache}, \texttt{Rede} und \texttt{Protokoll} erben von \texttt{Document}. Die Benutzerrollen erben von \texttt{HumanUsers}.
    \item \textbf{Realisierung:}  Eine Klasse implementiert ein Interface.
    \item \textbf{Aggregation:} Ein \texttt{Protokoll} enthält mehrere eigenständige \texttt{Rede}-Dokumente.
    \item \textbf{Assoziation:} Direkte Verbindung zwischen Klassen, z.\,B. \texttt{DocumentEditor} und \texttt{DocumentStorage}.
    \item \textbf{Abhängigkeit:}Eine Klasse verwendet eine andere, ohne sie als Attribut zu speichern.
\end{itemize}

\subsection{Kardinalitäten}
Nur bei Assoziationen und Aggregationen. Sie zeigen, wie viele Objekte miteinander verbunden sein können.

\begin{center}
\begin{tabular}{|p{4.2cm}|p{2.5cm}|p{8.5cm}|}
\hline
\textbf{Beziehung} & \textbf{Kardinalität} & \textbf{Bedeutung} \\ \hline
\texttt{Protokoll} -- \texttt{Rede} & 0..* zu 0..* & Ein Protokoll kann mehrere Reden enthalten. \\ \hline
\texttt{DocumentEditor} -- \texttt{DocumentStorage} & 0..* zu 1 & Jeder Editor hat genau einen Speicher. \\ \hline
\texttt{DocumentData} -- \texttt{DocumentType} & 0..* zu 1 & Jedes DocumentData hat genau einen Typ. \\ \hline
\texttt{DocumentJson} -- \texttt{DocumentType} & 0..* zu 1 & Jedes gespeicherte Dokument hat einen Typ. \\ \hline
\texttt{ChangeLogEntry} -- \texttt{ChangeType} & 0..* zu 1 & Jeder Log-Eintrag hat einen Änderungstyp. \\ \hline
\end{tabular}
\end{center}

\paragraph{Hinweis zu Logs.}
\texttt{ChangeLogEntry} speichert nur \texttt{documentId} und \texttt{userId}, keine direkten Objekte. Deshalb gibt es keine Assoziation zu \texttt{Document} oder \texttt{HumanUsers}.

\subsection{Prinzip der Architektur}
\begin{itemize}
    \item \texttt{Document}: Gemeinsame Eigenschaften aller Dokumente.
    \item \texttt{Drucksache}, \texttt{Rede}, \texttt{Protokoll}: Eigene zusätzliche Eigenschaften.
    \item Alle Attribute private; Änderungen nur über den \texttt{DocumentEditor}.
    \item \texttt{DocumentEditor}: Berechtigung prüfen, Änderungen durchführen, Log-Eintrag erstellen und speichern.
    \item \texttt{DocumentFactory}: Erstellt Dokumente anhand von \texttt{DocumentData}.
    \item \texttt{DocumentStorage}: Interface zum Speichern und Laden; \texttt{JsonDocumentStorage} nutzt JSON.
    \item \texttt{HumanUsers}: Benutzerrechte über \texttt{Berechtigung}.
    \item \texttt{ExternalAPIClient}: Zugriff nur auf öffentliche Dokumente.
\end{itemize}

\subsection{Nicht dargestellt}
Java-Klassen, JSON-Bibliotheken und Testklassen (z.\,B. \texttt{LocalDateTime}, \texttt{Path}, \texttt{ObjectMapper}).